\documentclass[
  customlatinfont = macos,
]{hkustthesis}

\usepackage{hkust-cjk}
\usepackage{colortbl}
\usepackage{tabularx}
\usepackage[most]{tcolorbox}

% Extra colors for tcolorbox highlights
\definecolor{moeBg}{HTML}{EEF2FF}
\definecolor{moeBorder}{HTML}{6366F1}
\definecolor{quantBg}{HTML}{F0FDF4}
\definecolor{quantBorder}{HTML}{22C55E}
\definecolor{longBg}{HTML}{FFF7ED}
\definecolor{longBorder}{HTML}{F97316}
\definecolor{ttcBg}{HTML}{FDF2F8}
\definecolor{ttcBorder}{HTML}{EC4899}
\definecolor{synthBg}{HTML}{F5F3FF}
\definecolor{synthBorder}{HTML}{8B5CF6}
\definecolor{mergeBg}{HTML}{ECFDF5}
\definecolor{mergeBorder}{HTML}{10B981}

\newtcolorbox{insightbox}[2][]{
  colback=#2Bg, colframe=#2Border,
  fonttitle=\bfseries, title={#1},
  boxrule=0.8pt, arc=3pt,
  before skip=10pt, after skip=10pt,
}

\hkustsetup{
  info = {
    degree        = {mphil},
    title         = {Open-Source LLM Core Techniques: A Deep Dive},
    author        = {Ziming Wang},
    school        = {Academy of Interdisciplinary Studies},
    department    = {Division of Integrative Systems and Design},
    program       = {Technology Innovation and Entrepreneurship},
    supervisor    = {},
    submit-month  = {March 2026},
    submit-date   = {7 March 2026},
    defend-date   = {},
    city          = {Hong Kong},
  }
}

\addbibresource{refs.bib}
\emergencystretch=2em

\begin{document}

\frontmatter
\tableofcontents

\mainmatter

%% ============================================================
\chapter{Mixture of Experts (MoE)}
\label{ch:moe}
%% ============================================================

MoE 架构已成为 2026 年所有 frontier 开源模型的标准选择。
其核心价值在于以极低的训练成本获得同等质量的模型：
DeepSeek-V3 以 \$5.6M 的成本训练了 671B 参数模型，
而西方同级别模型的训练成本高达数亿美元~\cite{deepseekv3}.

\section{工作原理}

对于每个输入 token，MoE 层的处理流程如下：
\begin{enumerate}[noitemsep]
  \item \textbf{路由}：一个学习到的线性门控网络计算每个专家的分数。
  \item \textbf{选择}：Top-K 选择激活 $K$ 个专家（通常 $K=2$ 或 $K=8$）。
  \item \textbf{并行计算}：被选中的专家独立处理 token。
  \item \textbf{加权聚合}：按路由权重对专家输出求和。
\end{enumerate}

相比于 dense 模型，MoE 使得总参数量远大于每个 token 实际激活的参数量。
例如 DeepSeek-V3 总参数 671B，但每个 token 仅激活 37B（$\sim$5.5\%）。

\section{模型格局（截至 2026 年 3 月）}

\begin{table}[H]
\centering\small
\caption{主要 MoE 模型对比}
\label{tab:moe-models}
\begin{tabularx}{\linewidth}{l r r r r r}
\toprule
\textbf{模型} & \textbf{总参数} & \textbf{活跃参数} & \textbf{专家数} & \textbf{Top-K} & \textbf{上下文} \\
\midrule
DeepSeek-V4       & $\sim$1T & 32B & 256+ & TBD & 1M \\
DeepSeek-V3       & 671B & 37B & 256+1 shared & 8 & 128K \\
Mistral Large 3   & 675B & 41B & MoE & --- & 256K \\
LLaMA 4 Maverick  & $\sim$400B & 17B & 128 & custom & 1M \\
Qwen3-235B        & 235B & 22B & 128 & 8 & 256K \\
GPT-OSS-120B      & 117B & 5.1B & MoE+MXFP4 & --- & --- \\
LLaMA 4 Scout     & 109B & 17B & 16 & custom & 10M \\
\bottomrule
\end{tabularx}
\end{table}

关键趋势：专家数量从 Mixtral 的 8 个猛增到 DeepSeek 的 256 个。
更细粒度的专家意味着更灵活的组合路由——
$\binom{256}{8} \approx 4.4 \times 10^{9}$ 种组合，
相比 $\binom{8}{2} = 28$ 种，组合空间扩大了 $10^{8}$ 倍。

\section{DeepSeek 三大创新}

\subsection{细粒度专家（Fine-Grained Experts）}

DeepSeek 使用 256 个路由专家加 1 个共享专家。
每个路由专家仅为标准 FFN 的 1/32 大小，
共享专家处理所有 token 以捕获通用知识：
\begin{equation}
  \mathbf{y} = \mathrm{SharedExpert}(\mathbf{x})
  + \sum_{i \in \mathrm{TopK}} g_i \cdot \mathrm{Expert}_i(\mathbf{x})
  \label{eq:moe-output}
\end{equation}
其中 $g_i$ 为路由权重，TopK 选取 8 个专家。

\subsection{Auxiliary-Loss-Free 负载均衡}

传统方法在训练目标中加入辅助损失 $L_{\mathrm{balance}}$
以惩罚专家利用率不均衡：
\begin{equation}
  L_{\mathrm{total}} = L_{\mathrm{LM}} + \alpha \cdot L_{\mathrm{balance}}
\end{equation}
这引入了根本矛盾：$\alpha$ 太大损害模型质量，太小则专家坍缩。

\begin{insightbox}[DeepSeek 的方案：Bias-Based Balancing]{moe}
不加辅助损失，而是给每个专家的路由分数加一个 \textbf{不参与反向传播} 的偏置项 $b_i$：
\[
  \mathrm{score}_i = \mathrm{softmax}\!\bigl(W_{\mathrm{router}} \cdot \mathbf{x} + b_i\bigr)
\]
偏置 $b_i$ 由简单启发式规则在每步更新：
专家过载则 $b_i \gets b_i - \gamma$，
空闲则 $b_i \gets b_i + \gamma$（$\gamma \approx 10^{-3}$）。
由于 $b_i$ 不在计算图内，语言建模梯度完全不受干扰。
2025 年芝加哥大学的理论论文证明该方法等价于
一步一迭代的原始--对偶分配问题求解，
具有单调改进保证和对数遗憾界。
\end{insightbox}

\subsection{DeepEP：专用 All-to-All 通信库}

MoE 需要跨 GPU 分发 token 到对应专家（all-to-all dispatch）。
DeepEP 提供两种模式：

\begin{table}[H]
\centering\small
\caption{DeepEP 两种内核模式}
\begin{tabularx}{\linewidth}{l l X}
\toprule
\textbf{模式} & \textbf{用途} & \textbf{技术特点} \\
\midrule
Normal（高吞吐） & 训练 + prefill & NVLink + RDMA 混合，SM 数量可控 \\
Low-Latency      & 推理 decode  & Hook-based 通信--计算重叠，不消耗 SM 资源 \\
\bottomrule
\end{tabularx}
\end{table}

\section{源码关键路径}

MoE 前向传播的完整流程可在以下开源项目中追踪：

\begin{table}[H]
\centering\small
\caption{各框架 MoE 核心代码文件}
\label{tab:moe-code}
\begin{tabularx}{\linewidth}{l l X}
\toprule
\textbf{框架} & \textbf{组件} & \textbf{关键路径} \\
\midrule
vLLM & 路由器 & \texttt{layers/fused\_moe/router/fused\_topk\_router.py} \\
vLLM & 融合内核 & \texttt{layers/fused\_moe/fused\_moe.py} (79KB, Triton) \\
vLLM & DeepEP 后端 & \texttt{deepep\_ht\_prepare\_finalize.py} \\
Megatron & 路由器 & \texttt{moe/router.py} (TopKRouter) \\
Megatron & Token 调度 & \texttt{moe/token\_dispatcher.py} \\
Megatron & 负载均衡损失 & \texttt{moe/moe\_utils.py} \\
DeepSpeed & 分片 MoE & \texttt{moe/sharded\_moe.py} \\
TorchTitan & EP 并行 & \texttt{distributed/expert\_parallel.py} (DTensor) \\
Nanotron & Grouped GEMM & \texttt{src/nanotron/nn/moe.py} \\
\bottomrule
\end{tabularx}
\end{table}


%% ============================================================
\chapter{Quantization（量化）}
\label{ch:quantization}
%% ============================================================

量化通过降低模型精度来减少内存占用、提高推理速度，
同时控制质量损失在可接受范围内。

\section{方法选择矩阵}

\begin{insightbox}[量化方法速查]{quant}
\begin{itemize}[noitemsep]
  \item \textbf{数据中心 GPU (H100/H200)}：FP8（vLLM/TRT-LLM），近零精度损失
  \item \textbf{Blackwell GPU (B200/RTX 5090)}：NVFP4，$2.3\times$ 吞吐提升
  \item \textbf{消费级 GPU + 高吞吐}：AWQ + Marlin kernel，741 tok/s（超过 FP16 的 461）
  \item \textbf{CPU / Apple Silicon}：GGUF Q4\_K\_M（llama.cpp / Ollama）
  \item \textbf{微调}：BitsAndBytes NF4（QLoRA）
\end{itemize}
\end{insightbox}

\section{H200 实测对比}

在 H200 GPU 上对 Qwen2.5-32B-Instruct 进行 4-bit 量化的实测结果如
\cref{tab:quant-bench} 所示。

\begin{table}[H]
\centering\small
\caption{H200 上 Qwen2.5-32B 各量化方法对比~\cite{jarvislabs-quant}}
\label{tab:quant-bench}
\begin{tabularx}{\linewidth}{l r r r}
\toprule
\textbf{方法} & \textbf{Perplexity} & \textbf{HumanEval} & \textbf{吞吐 (tok/s)} \\
\midrule
FP16 Baseline  & 6.56 & 56.1\% & 461 \\
Marlin-AWQ     & 6.84 & 51.8\% & \textbf{741} \\
BitsAndBytes   & \textbf{6.67} & 51.8\% & 168 \\
GGUF Q4\_K\_M  & 6.74 & 51.8\% & 93 \\
\bottomrule
\end{tabularx}
\end{table}

AWQ + Marlin 内核在 4-bit 下吞吐量反超 FP16（741 vs.\ 461 tok/s）。
这是因为更小的模型减少了内存带宽瓶颈——
量化不仅节省内存，还能提升推理速度。

\section{GGUF K-Quant 深度解析}

GGUF 的 K-quant 采用两级层级量化（hierarchical quantization）：

\begin{enumerate}[noitemsep]
  \item 层权重被切分为 \textbf{super-block}（如 256 个权重），
        配一个 FP16 缩放因子。
  \item 每个 super-block 内部包含多个 \textbf{sub-block}（如 32 个权重），
        各自独立计算缩放因子并量化权重。
  \item Sub-block 的缩放因子本身被\textbf{二次量化}（6-bit），
        由 super-block 级别的 FP16 缩放因子覆盖。
\end{enumerate}

\texttt{Q4\_K\_M} 命名含义：4-bit 精度 / K-quant 方法 / Medium 精度变体。
``M'' 表示混合精度策略：attention 层的 Q/K/V 和 FFN down\_proj 使用 Q6\_K，
其余层使用 Q4\_K。仅多 $\sim$0.2 bpw，perplexity 增量从 0.1149 降至 0.0535。

\begin{table}[H]
\centering\small
\caption{GGUF 量化格式对比（基于 LLaMA 7B）}
\label{tab:gguf}
\begin{tabularx}{\linewidth}{l r r r}
\toprule
\textbf{格式} & \textbf{有效 bpw} & \textbf{7B 大小} & \textbf{PPL $\Delta$} \\
\midrule
Q8\_0 & 8.0 & 6.7\,GB & +0.0004 \\
Q6\_K & 6.0 & 5.15\,GB & +0.0044 \\
Q5\_K\_M & 5.1 & 4.45\,GB & +0.0142 \\
\rowcolor{quantBg}
\textbf{Q4\_K\_M} & \textbf{4.5} & \textbf{3.80\,GB} & \textbf{+0.0535} \\
Q3\_K\_M & 3.3 & 3.06\,GB & +0.2437 \\
Q2\_K & 2.5 & 2.67\,GB & +0.8698 \\
\bottomrule
\end{tabularx}
\end{table}

\section{FP8 训练时量化}

DeepSeek-V3 是首个 FP8 端到端训练的大规模模型~\cite{deepseekv3}。核心技术包括：

\begin{enumerate}[noitemsep]
  \item \textbf{细粒度块级量化}：权重每 $128 \times 128$ 子矩阵独立缩放，
        激活每 $1 \times 128$ 子向量独立缩放。
  \item \textbf{高精度累加}：FP8 Tensor Core 仅提供 14-bit 精度，
        DeepSeek 通过 4 次连续 WGMMA 操作后加入 FP32 累加器解决此问题。
  \item \textbf{PyTorch 原生支持}：\texttt{torchao.float8} 的 rowwise recipe
        在 2K H200 集群上实现 1.34--1.43$\times$ 训练加速。
\end{enumerate}


%% ============================================================
\chapter{Long Context（长上下文）}
\label{ch:long-context}
%% ============================================================

\section{上下文长度排行}

\begin{table}[H]
\centering\small
\caption{当前最长上下文窗口模型}
\label{tab:long-ctx}
\begin{tabularx}{\linewidth}{l r X}
\toprule
\textbf{模型} & \textbf{上下文} & \textbf{关键技术} \\
\midrule
LLaMA 4 Scout    & 10M  & iRoPE（RoPE + NoPE 交替层） \\
Grok-4-fast       & 2M   & 未公开 \\
Gemini 3 Pro      & 1M   & Sparse MoE \\
Claude 4.5 Sonnet & 1M   & 未公开 \\
DeepSeek-V4       & 1M   & MLA + DSA \\
Qwen2.5-1M        & 1M   & YaRN + 渐进式预训练 \\
\bottomrule
\end{tabularx}
\end{table}

\section{核心技术演进}

\subsection{YaRN（Yet Another RoPE extensioN）}

RoPE 在训练长度之外会"外推失败"，位置编码进入未见过的频率空间。
YaRN 整合了三个关键创新：

\begin{enumerate}[noitemsep]
  \item \textbf{NTK-By-Parts}：将频率维度分为三组——
        低频不动（保留原始模式）、中频 NTK-aware 缩放、高频激进缩放。
  \item \textbf{Temperature Scaling}：
        在 softmax 中加入温度参数，短序列保持原始行为，长序列时稳定注意力分布。
  \item \textbf{最小微调}：仅需原始预训练数据的极小比例即可实现巨大的上下文扩展。
\end{enumerate}

目前 Qwen、DeepSeek、LLaMA 等主流模型均使用 YaRN 变体进行上下文扩展。

\subsection{iRoPE（Meta，LLaMA 4）}

LLaMA 4 交替使用两种注意力层：
\begin{itemize}[noitemsep]
  \item \textbf{RoPE 层}：带位置编码，捕捉局部上下文和相对位置。
  \item \textbf{NoPE 层}：无位置偏置，关注纯语义关系。
\end{itemize}

NoPE 层天然不受序列长度限制，
而 RoPE 层通过 temperature scaling 在推理时泛化到超出训练长度，
共同支撑了 LLaMA 4 Scout 的 10M token 上下文窗口。

\subsection{Ring Attention}

将长序列分割到多个 GPU 上，设备排列成环形拓扑，
KV block 在设备间循环传递。
每个设备计算自己持有的 query 与当前可见 KV 的注意力，
通信与计算重叠以隐藏传输延迟。
内存复杂度从 $O(n^2)$ 降为每设备 $O(n)$。

\section{Context Rot：最重要的研究发现}

\begin{insightbox}[Chroma 2025 年实证研究：18 个 frontier 模型]{long}
\textbf{所有模型都随输入长度增加而性能下降，即使远未达到上下文窗口上限。}
\begin{itemize}[noitemsep]
  \item \textbf{U 型曲线}：模型对序列开头和结尾的信息关注度最高，
        中间信息的检索准确度下降超过 30\%。
  \item \textbf{反直觉发现}：有结构的上下文比打乱的上下文更难检索——
        模型在有逻辑的干扰信息中更容易"迷路"。
  \item \textbf{核心结论}：Context engineering（如何组织上下文）
        比 context length（能塞多少）更重要。
\end{itemize}
\end{insightbox}

\section{推理框架的长上下文支持}

\textbf{vLLM} 提供 PagedAttention（KV cache 分页管理）、
Continuous Batching、Chunked Prefill 和 Prefix Caching。

\textbf{SGLang} 的 Chunked Pipeline Parallelism（2026 年 1 月发布）
在 DeepSeek-V3.1 上实测实现 $3.31\times$ prefill 吞吐提升，
Qwen3-235B 处理 1M token 的 TTFT 从 55.5s 降至 10.5s（降低 81\%）。


%% ============================================================
\chapter{Test-Time Compute（推理时计算缩放）}
\label{ch:ttc}
%% ============================================================

\section{核心洞察}

\begin{insightbox}[推理时计算的颠覆性影响]{ttc}
\textbf{7B 模型 + 100$\times$ 推理计算 $\approx$ 70B 模型标准推理。}
推理需求预计 2026 年超过训练需求 118 倍。
GPT-4 等效性能成本从 2022 年 \$20/M tokens 降至 2026 年 \$0.40/M tokens
（$1000\times$ 下降）。
\end{insightbox}

\section{三种策略}

\begin{table}[H]
\centering\small
\caption{推理时计算缩放的三种策略}
\begin{tabularx}{\linewidth}{l l X}
\toprule
\textbf{策略} & \textbf{方法} & \textbf{代表} \\
\midrule
Parallel   & 采样 $N$ 个答案取共识 & Best-of-N, Majority Vote \\
Sequential & 链式推理，逐步验证   & DeepSeek-R1 thinking tokens \\
Hybrid     & 并行 + 顺序结合     & ThreadWeaver ($1.53\times$ 延迟降低) \\
\bottomrule
\end{tabularx}
\end{table}

\section{推理模型格局}

\begin{table}[H]
\centering\small
\caption{主要推理模型对比（2026 年初）}
\label{tab:reasoning-models}
\begin{tabularx}{\linewidth}{l X}
\toprule
\textbf{模型} & \textbf{特点} \\
\midrule
DeepSeek-R1-0528 & 671B MoE, AIME 87.5\%, 开源 MIT \\
Qwen3            & 统一 thinking/non-thinking 模式，单模型两种用法 \\
o3/o4-mini       & 首个推理 + 全工具访问，ARC-AGI 87\%+ \\
Claude Opus 4.6  & 可配置 thinking budget (32K/64K tokens) \\
\bottomrule
\end{tabularx}
\end{table}

\section{GRPO vs.\ REINFORCE++ vs.\ PPO}

这三种算法是训练推理模型的核心 RL 方法。

\textbf{GRPO}（Group Relative Policy Optimization）的核心思想是
消除 PPO 中的 critic network，使用组内采样估算 advantage：
\begin{equation}
  A_i = \frac{R_i - \mathrm{mean}(\{R_j\})}{\mathrm{std}(\{R_j\})}
  \label{eq:grpo-advantage}
\end{equation}
\begin{equation}
  J_{\mathrm{GRPO}} = \mathbb{E}\!\left[
    \min\!\left(\frac{\pi_\theta}{\pi_{\mathrm{old}}} A_i,\;
    \mathrm{clip}\!\left(\frac{\pi_\theta}{\pi_{\mathrm{old}}},\,1\!\pm\!\epsilon\right) A_i\right)
  \right] - \beta \cdot D_{\mathrm{KL}}(\pi_\theta \| \pi_{\mathrm{ref}})
  \label{eq:grpo-objective}
\end{equation}

\begin{table}[H]
\centering\small
\caption{三种 RL 算法对比}
\label{tab:rl-compare}
\begin{tabularx}{\linewidth}{l X X X}
\toprule
\textbf{维度} & \textbf{PPO} & \textbf{GRPO} & \textbf{REINFORCE++} \\
\midrule
Advantage 估计 & Critic network $V(s)$ & 组内 reward 均值/标准差 & 组内均值 + batch-level 标准差 \\
内存需求 & 4 份模型副本 & 3 份（无 critic），内存减半 & 同 GRPO \\
稳定性 & 最稳定 & 存在 reward collapse & 更鲁棒 \\
可扩展性 & 70B+ 极难 & 可行 & 可行 \\
\bottomrule
\end{tabularx}
\end{table}

2026 年共识：GRPO 最流行（DeepSeek-R1 验证），
REINFORCE++ 更稳定（逐步替代），
PPO 在 70B+ 因内存限制几乎不可行。

\section{DeepSeek-R1 蒸馏管线}

DeepSeek-R1 的完整训练经历五个阶段：
\begin{enumerate}[noitemsep]
  \item \textbf{Base}：DeepSeek-V3-Base（671B MoE 预训练模型）。
  \item \textbf{纯 RL}：使用 GRPO 训练 $\to$ R1-Zero。
        自然涌现 chain-of-thought、自我验证、aha moments。
  \item \textbf{Cold-start SFT + 二轮 RL}：先用精选数据做 SFT，再做第二轮 RL $\to$ R1。
  \item \textbf{蒸馏}：R1 生成 $\sim$800K 推理 trace，
        SFT 训练小模型 $\to$ R1-Distill-Qwen-1.5B/7B/14B/32B,
        R1-Distill-Llama-8B/70B。
  \item \textbf{对齐}：安全对齐和行为稳定。
\end{enumerate}

开源复现项目 HuggingFace Open-R1 (19.9K stars)
完整实现了上述三阶段方法。


%% ============================================================
\chapter{Synthetic Data（合成数据）}
\label{ch:synthetic}
%% ============================================================

合成数据已成为 frontier 模型训练的标准组件。
2026 年的共识策略是\textbf{合成 + 精选人类数据的混合方法}。

\section{前沿模型的合成数据策略}

\begin{table}[H]
\centering\small
\caption{前沿模型的合成数据使用方式}
\begin{tabularx}{\linewidth}{l X}
\toprule
\textbf{模型} & \textbf{策略} \\
\midrule
Phi-4 (Microsoft) & 仅 140 万精选样本 SFT，证明数据质量 $>$ 数据规模 \\
Qwen3             & 36T tokens 含大量用 Qwen2.5-Math/Coder 生成的领域合成数据 \\
DeepSeek-R1       & 纯 RL 训练推理能力，再用 R1 推理 trace 蒸馏到小模型 \\
\bottomrule
\end{tabularx}
\end{table}

\section{最优合成--真实数据比例}

Meta FAIR 训练超过 1000 个 LLM 的大规模研究得出核心结论：
\begin{insightbox}[Meta FAIR 研究结论]{synth}
预训练最优比例：$\sim$30\% 合成 + $\sim$70\% 真实数据，可加速训练 5--10$\times$。
\begin{itemize}[noitemsep]
  \item 单轮 rephrased 合成数据不会导致 model collapse。
  \item 递归生成和纯 textbook 风格数据会导致 collapse。
  \item 8B 参数的生成器产出效果不亚于更大模型。
\end{itemize}
\end{insightbox}

\section{Model Collapse 风险}

在合成数据上递归训练会导致模型"遗忘"数据分布的尾部，
逐代性能退化。缓解策略是始终混入精选的人类/网络数据，
典型比例为 50--80\% 合成 + 20--50\% 人类数据。


%% ============================================================
\chapter{Model Merging（模型合并）}
\label{ch:merging}
%% ============================================================

模型合并将多个微调后的模型组合成一个，
无需额外训练。仅用 CPU（无需 GPU）即可通过 MergeKit 实现。

\section{核心方法}

\begin{table}[H]
\centering\small
\caption{模型合并方法对比}
\begin{tabularx}{\linewidth}{l l X}
\toprule
\textbf{方法} & \textbf{原理} & \textbf{备注} \\
\midrule
SLERP          & 球面线性插值 & 两个模型平滑混合 \\
TIES-Merging   & 修剪 + 符号选举 + 合并 & 处理参数冲突 \\
DARE           & 丢弃 90--99\% task vector 权重 & 极致稀疏化 \\
Task Arithmetic & 加减 task vector & 简单而有效 \\
进化式合并      & 进化算法搜索最优配置 & 2025 年最大突破 (Sakana AI) \\
\bottomrule
\end{tabularx}
\end{table}

2025 年的最大突破来自 Sakana AI 的\textbf{进化式模型合并}，
发表于 Nature Machine Intelligence，
使用进化算法搜索最优合并配置（层级、权重系数），
已集成到 MergeKit 中。

MergeKit（Arcee AI，6.8K stars）是该领域无可争议的标准工具，
CPU 即可运行。NeurIPS 2024 举办了正式的 LLM Merging Competition，
标志着模型合并已从实验技术进入主流。


%% ============================================================
\chapter{推荐学习路径与资源}
\label{ch:resources}
%% ============================================================

\begin{table}[H]
\centering\small
\caption{按主题推荐的核心阅读资源}
\label{tab:resources}
\begin{tabularx}{\linewidth}{l l X}
\toprule
\textbf{主题} & \textbf{优先级} & \textbf{资源} \\
\midrule
MoE & 必读 & Cameron Wolfe: nanoMoE (手写 PyTorch 实现) \\
MoE & 必读 & Sebastian Raschka: Big LLM Architecture Comparison \\
MoE & 推荐 & Shirley Li: DeepSeek-V3 Explained (8 篇系列) \\
\midrule
量化 & 必读 & JarvisLabs: vLLM 量化完全指南 (H200 实测) \\
量化 & 必读 & Colfax Research: DeepSeek FP8 训练 \\
量化 & 推荐 & llama.cpp PR \#1684 (K-quant 原始实现) \\
\midrule
长上下文 & 必读 & Chroma: Context Rot (18 模型实证) \\
长上下文 & 必读 & Aman Arora: How LLMs Scaled from 512 to 2M Context \\
长上下文 & 推荐 & LMSYS: SGLang Pipeline Parallelism \\
\midrule
推理计算 & 必读 & Lilian Weng: Why We Think \\
推理计算 & 必读 & The Illustrated GRPO \\
推理计算 & 推荐 & Cameron Wolfe: GRPO++ \\
\midrule
合成数据 & 推荐 & Meta FAIR: 1000+ LLM 合成数据实验 \\
模型合并 & 推荐 & Maxime Labonne: Merge Models with MergeKit \\
\bottomrule
\end{tabularx}
\end{table}


\printbibliography[heading=bibintoc,title=References]

\end{document}
